% Keksitty graafisen ohjeiston kuva: Organisaatio Oy:n logon suoja-alue.
% Suoja-alueen leveys x on logon merkkiosan korkeus. Rajataan
% preview-paketilla; logo-suoja-alue.pdf on esimerkki-kayttoohje.tex-
% asiakirjan käyttämä kuva.
\documentclass{article}
\usepackage[active,tightpage]{preview}
\usepackage{helvet}
\usepackage{tikz}
\usetikzlibrary{calc}
\setlength{\PreviewBorder}{1pt}
\PreviewEnvironment{tikzpicture}

\begin{document}
\begin{tikzpicture}
  % Logo (sama piirros kuin logo-organisaatio.tex:ssä)
  \begin{scope}[local bounding box=logo]
    \fill[cyan!75!blue,opacity=.8]  (0,0)     circle (3.4mm);
    \fill[blue!65!cyan,opacity=.8]  (3.4mm,0) circle (3.4mm);
    \fill[blue!55!black,opacity=.8] (6.8mm,0) circle (3.4mm);
    \node[anchor=west,text=blue!50!black,
          font=\sffamily\bfseries\large] at (11.6mm,0) {Organisaatio Oy};
  \end{scope}
  % Suoja-alueen raja: x = merkin korkeus 6,8 mm joka suuntaan
  \draw[gray,dashed]
    ($(logo.south west)+(-6.8mm,-6.8mm)$) rectangle
    ($(logo.north east)+(6.8mm,6.8mm)$);
  % Mittanuolet
  \draw[gray,<->] ($(logo.west)+(-6.8mm,0)$) -- (logo.west)
    node[midway,above,font=\sffamily\scriptsize,text=gray] {x};
  \draw[gray,<->] (logo.north) -- ($(logo.north)+(0,6.8mm)$)
    node[midway,right,font=\sffamily\scriptsize,text=gray] {x};
  % Mitan määritelmä: x = merkin korkeus
  \draw[gray,<->] ($(logo.south east)+(2.5mm,0)$) -- ($(logo.north east)+(2.5mm,0)$)
    node[midway,right,font=\sffamily\scriptsize,text=gray] {x};
\end{tikzpicture}
\end{document}
